\section{Amplification along observations}
\label{prof:amplification}

\subsection{A compact-interval improvement of overlap tails}
\label{prof:sec-bridge}

We next improve the probability of the static covariance bounds.  The first
ingredient concerns overlaps about the exact field root on a compact interval
of small positive observation times.  Throughout this section,
\(\mathcal G\) denotes the spectral restriction already imposed,
\(r=\sqrt s\), and \(k=nr^3\).  Probabilities with the indicator of
\(\mathcal G\) are integrated probabilities under the planted law; no claim
about the planted probability of \(\mathcal G\) is needed.  We write
\(\Sigma_s=\operatorname{Cov}_{\mu_{h_s}}X\) and
\(m_s^{\mathrm{post}}=\mu_{h_s}(X)\).  Thus the subscripted vector
\(m_s^{\mathrm{post}}\) is distinct from the posterior probability measure
\(\mu_{h_s}\), and from the exact root \(m\).

\begin{lemma}[Overlap improvement on compact intervals]
\label{prof:lem-compact-bridge}
There are \(r_0>0\) and \(a_*>0\), with \(a_*\) independent of a lower
endpoint \(r_1\), such that the following holds.  For every fixed
\(0<r_1\le r_0\) and \(\chi>0\), there is \(c=c(r_1,r_0,\chi)>0\)
for which, at each \(r\in[r_1,r_0]\),
\[
 \mu_{h_s}^{\otimes2}
 \bigl( |(X-m)^T(X'-m)|>\chi nr\bigr)\le e^{-cn}
\]
outside an integrated planted exceptional set of mass at most
\(e^{-a_*k}\), with \(\mathcal G\) imposed.  The root is the unique
regular root on the retained event.
\end{lemma}

\begin{proof}
Use the root orbits, their three projection columns \(P\), and the reference
partition integral \(Z_*\) from the local comparison.  We write
\(\vartheta\) for the normalized projection coordinate, so that
\(\delta=\operatorname{diag}(\sqrt q,q,q^{3/2})\vartheta\), and put
\(z=X-m\), \(N=nq^3\).  For different replicas let
\[
 t_{ij}=n^{-1}(z_i-P\delta_i)^T(z_j-P\delta_j).
\]
On this compact interval all polynomial bounds may depend on \(r_1\).

First, outside Haar mass \(e^{-Bk}\), with any prescribed fixed \(B\),
the full tilted spin integral is at least \(Z_*e^{-o(n)}\).
Indeed \eqref{prof:mass-anchor} supplies a fixed positive fraction of
orientations with full mass at least \(Z_*/2\). The simultaneous
covariance and mean bounds in \eqref{prof:eq13} give a bounded raw
operator second moment of \(z\) on this compact interval.
Jensen's inequality, evaluated at such an anchor, bounds the
loss under a rotation at fixed Frobenius distance by
\(C\sqrt n+O(1)\): the linear energy increment pairs with a raw
second-moment matrix of Hilbert--Schmidt norm \(O(\sqrt n)\), and the
quadratic increment costs \(O(1)\).

Concentration of distance to the anchor
set excludes distances greater than a sufficiently large fixed constant
with probability \(e^{-Bk}\).  The determinant weight changes this estimate
only by the already established arbitrarily small loss in its rate.

Fix a small absolute \(D_*>0\).  For any fixed replica number \(l\),
consider the unnormalized tilted spin integral on
\begin{equation}
 D_*\le \sum_{i=1}^l|\vartheta_i|^2
       +\sum_{i<j}|t_{ij}/q|^2\le4D_*.
 \label{prof:eq-bridge-annulus}
\end{equation}
Its Haar mean, divided by \(Z_*^l\), is at most
\(e^{-cD_*N}\), where \(c>0\) is independent of \(l\), although the
threshold in \(n\) may depend on \(l\).  We verify this uniformity because
the replica number will depend on \(\chi\).

Under the tilted reference, the projection coordinates
\(\sqrt N\,\vartheta_i\) have bounded Gaussian covariance before length
conditioning.  Every linear combination with unit Euclidean coefficients
of the off-diagonal coordinates \(nt_{ij}/\sqrt{n/q}\) is a centered
Gaussian quadratic form whose matrix has bounded Hilbert--Schmidt norm and
operator norm at most \(C/\sqrt N\).  These constants are independent of
\(l\).  Gaussian exponential Markov bounds, followed by a fixed-mesh net,
therefore give cost at least \(c_0D_*N\) for
\eqref{prof:eq-bridge-annulus}.  The net and conditioning on the finitely
many lengths cost only polynomial factors for fixed \(l\); the latter
follows either from the length-density bounds or from buffered shells.

The annulus penalty just obtained is under the \emph{physically
tilted} reference shell, with covariance \(C_o\) before length
conditioning. The change of prior uses a different density:
the \emph{unweighted} \(o\) prior has uniform complement directions
after its projections and lengths are fixed. The common angular energy
integral is multiplied by the spin Gram density divided by this reference
Gram density. We estimate this quotient using polynomial boxes and
Chernoff upper bounds.
The single-projection density ratios in~\eqref{prof:eq14} cost at most
\(CN D_*^{3/2}+o(n)\).  At fixed projections the unweighted reference conditional
Gram density has exponent
\[
 -\frac n2\sum_{i<j}\frac{t_{ij}^2}{B_iB_j}
       +O(ND_*^{3/2}),
 \qquad
 B_i=b-2\sqrt q\,(\delta_i)_1-|\delta_i|^2,
\]
up to polynomial factors.  At the independent spin projection saddles
write the site deviations as \(\xi_i+d_i^0\), with independent centered
bounded \(\xi_i\).

Set \(R_{ij}=t_{ij}/(B_iB_j)\) and \(R_{ii}=0\).
The quadratic log moment satisfies
\[
 \log\mathbb E\exp(\xi^TR\xi/2)
  =\tfrac12\operatorname{Var}(\xi^TR\xi/2)
       +O(\|R\|_{\mathrm{HS}}^3),
\]
with an absolute constant when \(\|R\|_{\mathrm{HS}}\) is small.
For completeness, decouple the off-diagonal quadratic form using a random
bipartition.  Jensen's inequality bounds its exponential moment by that
of four times the corresponding directed cross sum.  Integrating one
side by the bounded-variable Gaussian moment bound produces a squared
linear-image penalty; introducing a white Gaussian and integrating the
other side gives a uniform exponential moment for a small multiple of
\(|\xi^TR\xi|/\|R\|_{\mathrm{HS}}\).  Taylor's formula then gives the
displayed third-order remainder.

The linear terms involving \(d_i^0\) are treated by H\"older's inequality
with relative inflation \(O(q\sqrt{D_*})\) in the centered quadratic
moment.  The remaining averaged log cost is at most
\[
 C\bigl(1+(q\sqrt{D_*})^{-1}\|R\|\bigr)
       \|R\|q^3D_*.
\]
Here we used the response bounds~\eqref{prof:eq17} for \(d_i^0\) and linear sub-Gaussianity.
The site-averaged quadratic coefficient for each pair differs from
\(B_iB_j\) by \(O(q^2)\). To see that the constant is independent
of \(l\), note on the annulus that
\[
 \|R\|_{\rm HS}\le Cq\sqrt{D_*},\qquad
 \frac1n\sum_{\text{sites}}\sum_{i=1}^l|d_i^0|^2
 \le Cq^3D_*.
\]
The latter inequality is \eqref{prof:eq17} summed over replicas.
The displayed log-moment bound is first applied at each site and
then averaged. Its cubic remainder costs
\(n\|R\|_{\rm HS}^3\le CN D_*^{3/2}\), the variance-profile
discrepancy costs
\(nq^2\|R\|_{\rm HS}^2\le CNqD_*\), and the displayed
H\"older cost, after multiplication by \(n\), is at most
\(CNqD_*^{3/2}\). These estimates use only sums of squares and
the Hilbert--Schmidt norm. Their constants therefore do not grow
with the fixed replica number. After the Chernoff subtraction, the
total relative cost is bounded by
\(CN D_*(q+\sqrt{D_*})\), besides polynomial factors. Choosing first
\(D_*\), then the upper bound on \(q\), absorbs this cost into
\(c_0D_*N\).  Angular integrals on these boxes have controlled variation;
their constants may depend on \(l,r_1\).

The annulus rate is independent of the fixed replica count. This lets
us turn a requested overlap accuracy into a sufficiently large, but
still fixed, number of replicas. A first-crossing selection will then
find an annulus subset without allowing a single large statistic to
account for the whole crossing.

We need an exponential posterior bound for each individually large
statistic, not just the polynomial cutoff loss used in the local
comparison. Let \(\mathcal R_s\) be the domain where all root/scalar
diagnostics retained for the local comparison hold. On this domain the
root is unique and regular, \(q>0\), and the three axes are nondegenerate.
Its excluded complement, including nonunique-root fields, has separate charge
\[
 \Pr_{\rm pl}(\mathcal G\cap\mathcal R_s^c)\le n^Ce^{-a_0k},
\]
after taking the fixed \(a_0>0\) below the preceding exclusion rates.

Only on \(\mathcal R_s\), define
\(E_1=\{|\vartheta_1|^2>D_*/4\}\) for one fresh posterior replica and
\(E_2=\{|t_{12}/q|^2>D_*/4\}\) for two. The plant and retained
root/scalar data remain fixed inside \(\mu_{h_s}^{\otimes j}\).
The scalar-root estimates after equivariant plant replacement, and
the plant--replica and replica--replica overlap estimates, give
\[
 \int_{\mathcal G\cap\mathcal R_s}
       \mu_{h_s}^{\otimes j}(E_j)\,d\Pr_{\rm pl}
 \le n^Ce^{-a_0k},\qquad j=1,2,
\]
after decreasing \(a_0\) if necessary. Conditional Markov on this
retained domain gives
\[
 \Pr_{\rm pl}\!\left\{\omega\in\mathcal G\cap\mathcal R_s:
       \mu_{h_s}^{\otimes j}(E_j)>e^{-a_0k/2}\right\}
 \le n^Ce^{-a_0k/2}.
\]
Relabeling \(E_1,E_2\) and taking a finite union covers every
individual statistic in any fixed replica family. The complement of
\(\mathcal R_s\) remains the separate charge displayed above.

Suppose the conclusion of the lemma fails. Draw a sufficiently large
fixed number \(l_{\rm pair}\) of disjoint violating pairs, depending
on \(\chi\). It uses \(2l_{\rm pair}\) replicas; the preceding estimates
apply to every subset with \(l\) equal to that subset's replica count. Since
\[
 n^{-1}z_i^Tz_j=t_{ij}+\delta_i\cdot\delta_j,
\]
their statistics force the total in \eqref{prof:eq-bridge-annulus} to cross
\(D_*\). Requiring all \(l_{\rm pair}\) pairs costs at most
\(l_{\rm pair}c n\) in log likelihood. Choose \(c\) so that
\(l_{\rm pair}c n\le a_0k/4\), possible uniformly on the compact
interval because \(k\ge nr_1^3\).
The preceding exponential conditional bounds then make all
individually large statistics, including those across pairs,
negligible even after this conditioning. This choice explains the
dependence of \(c\) on \(r_1,\chi\), while the field-rate constant
remains fixed.

A subset of the selected replicas then satisfies
\eqref{prof:eq-bridge-annulus}.  Indeed take the first prefix crossing
\(D_*\).  If its new row contributes too much, retain enough earlier
indices to make that row contribution comparable to \(D_*\); each of its
terms is at most \(D_*/4\), while the earlier internal sum is below
\(D_*\).  A finite union over subsets, the denominator lower bound, and
the preceding Haar estimate exclude the violating orientations at a rate
bounded below independently of the number of replicas.  Returning to the
counting weight costs \(o(k)\) on this compact interval.  The initially
excluded orbits had a fixed small rate as well.  This proves the lemma.
\end{proof}

\subsection{A conditional weak Poincar\'e inequality}
\label{prof:sec-conditional-variance}

The next lemma transports variance one short step backwards in the
observation filtration.  This Gaussian posterior interpolation belongs
to the stochastic-localization framework of
\citet{Eldan2013StochasticLocalization,ChenEldan2025Localization};
the exceptional-set variance estimate needed here is proved below.
Given \((J,h_s)\), let \(u=(1+d)s\),
\(\tau=u-s\), with \(d>0\) fixed and small, and write
\[
 g=\frac{h_u-h_s-\tau m_{\mathrm{sh}}}{\sqrt\tau}.
\]
The shift \(m_{\mathrm{sh}}\) is measurable at time \(s\) and has
polynomially bounded norm.  Let \(P\) be the conditional law of \(g\),
\(G\) standard Gaussian measure, and \(L_P=\log(dP/dG)\).
The Gaussian-mixture formula gives convexity of \(L_P\) and
\[
 \nabla^2\log\frac{dP}{dg}=-I+\tau\Sigma_u,
\]
where \(dP/dg\) is the Lebesgue density.

\begin{lemma}[Conditional variance estimate]
\label{prof:lem-conditional-variance}
Fix an observation bound \(T<\infty\) and positive constants
\(\alpha,\beta\). Suppose \(0<s<u=(1+d)s\le T\),
\(k=ns^{3/2}\), and
\begin{equation}
 P\{\tau\|\Sigma_u\|>1/4\}\le e^{-\beta k},
 \qquad
 \log\int e^{2L_P}\,dG\le\alpha k,
 \label{prof:eq20}
\end{equation}
where \(\alpha/\beta\) is less than a sufficiently small absolute
constant.  For every fixed \(D>0\) and smooth test \(\phi\) whose values
and required derivatives have fixed polynomial bounds,
\begin{equation}
 \operatorname{Var}_P\phi
 \le3\mathbb E_P\|\nabla_g\phi\|^2+n^{-D},
 \label{prof:eq21}
\end{equation}
provided \(\beta k/\log n\) is sufficiently large in terms of those
bounds and \(D\).  Bounded tests with polynomial cutoffs are included.
\end{lemma}

\begin{proof}
We first transfer small exceptional sets from \(P\) to \(G\).  Since
\(\int e^{L_P}dG=1\), the second moment in \eqref{prof:eq20} and
Paley--Zygmund give
\[
 G\{L_P\ge-1\}\ge(1-e^{-1})^2e^{-\alpha k}
                       \ge e^{-2\alpha k}
\]
for sufficiently large \(\alpha k\); enlarging harmless constants covers
the other case.  Also
\(G\{L_P>\beta k/8\}\le e^{-\beta k/8}\).
Convexity implies that
\begin{equation}
 G\{L_P<-\beta k/4\}\le e^{-c\beta k}
 \label{prof:eq-likelihood-lower-tail}
\end{equation}
for an absolute \(c>0\).

Here is a proof that does not presume regular
boundaries of the two level sets.  Use Gaussian endpoint conditioning to
drive two Brownian motions with the same noise into
\(\{L_P\ge-1\}\) and a proposed low set of mass at least
\(e^{-c\beta k}\), with joint success probability greater than \(0.8\).
Positive smooth approximations to the conditioning densities and energy
stopping, in the entropy--drift representation of
\citet[Proposition~1 and Section~2.5]{Lehec2013Entropy}, give total
expected squared drift energy at most
\(C(\alpha+c\beta)k\).

Reflect the second control about the first.
The third endpoint is the reflected point, so convexity forces its
likelihood above \(\beta k/8\) on the joint hit, for large \(k\).
Its relative entropy with respect to \(G\) is at most
\(C'(\alpha+c\beta)k\), by Girsanov and data processing.  Binary entropy
applied to \(\{L_P>\beta k/8\}\) contradicts this bound if \(c\) and
\(\alpha/\beta\) are small enough.  This proves
\eqref{prof:eq-likelihood-lower-tail}.  Consequently, if
\(P(A)\le e^{-\beta k/2}\), then
\[
 G(A)\le e^{-\beta k/4}+e^{-c\beta k}.
\]

The likelihood comparison has transferred rare starting sets to the
Gaussian reference. We now obtain contraction along typical Langevin
paths and telescope short-time variances to get the gradient term.
A final comparison through the Gaussian reference will remove the
remaining terminal variance.

Run stationary Langevin diffusion for \(P\), with generator
\(\Delta+\nabla\log(dP/dg)\cdot\nabla\), for
\(T'=C_D\log n\).  Apart from probability
\(\operatorname{poly}(n)e^{-\beta k}+e^{-n^2}\), its log-density Hessian
is at most \(-2I/3\) all along the path. Here are the buffers needed
for this path statement. Let \(M_n\) be a polynomial bound for the
Lipschitz constant of the Hessian, supplied by the bounded-spin
formulas for posterior derivatives, and take
\(\rho=(100M_n)^{-1}\). At every point of a polynomial time grid,
stationarity and \eqref{prof:eq20} give the upper bound \(-3I/4\)
outside total probability \(\operatorname{poly}(n)e^{-\beta k}\).
The drift has linear growth with a polynomial additive bound.
Choose the grid mesh polynomially small enough that, on a
polynomial ball, its drift and Brownian increments are less than
\(\rho/4\) between successive grid points. Gaussian increment
bounds and the tails of the Gaussian mixture \(P\) make the
remaining probability at most \(e^{-n^2}\). The Hessian Lipschitz
bound then preserves a strict upper bound below \(-2I/3\)
throughout the tube of radius \(\rho\) about the path.
The fixed observation bound and fixed \(\beta\) give
\(\beta k=O(n)\). Thus these Gaussian-tail errors remain negligible
after the exponential reweightings below. The lemma is only asserted
on this compact observation domain.

This tube also controls neighboring flows for that same future
noise. Start a synchronous flow within \(\rho/4\) of the path's
initial point and stop at the first time its distance from the
path reaches \(\rho/2\). Before that time, the mean Hessian along
the line segment between the flows is at most \(-2I/3\), so their
distance contracts. It cannot reach the stopping boundary.
Differentiating this stopped comparison and keeping a little slack
gives Jacobian norm at most \(e^{-v/2}\) for every \(v\le T'\).
This is a neighborhood of a good random path, not a claim about
every initial point.

Let \(K_v\) be the Langevin semigroup. Choose a polynomially small
\(\Delta t\) and an integer \(N_t\) with \(T'=N_t\Delta t\),
adjusting \(T'\) negligibly. Stationarity gives the exact telescope
\[
 \operatorname{Var}_P\phi-\operatorname{Var}_P K_{T'}\phi
 =\sum_{j=0}^{N_t-1}\E_{x\sim P}
   \operatorname{Var}_{K_{\Delta t}(x,\cdot)}
                  (K_{j\Delta t}\phi).
\]
Every short transition \(K_{\Delta t}(x,\cdot)\) has Poincar\'e constant
at most \((2+o(1))\Delta t\), uniformly in \(x\).  Indeed the global
upper Hessian bound is polynomial, hence the flow Jacobian is
\(1+o(1)\) up to this time; the conditional heat variance identity
integrated over the short interval gives the assertion.

For a variance term involving \(K_v\phi\), Jensen in the common future
noise bounds
\(\operatorname{Var}_{K_{\Delta t}(x,\cdot)}K_v\phi\) by the average
variance of \(y\mapsto\phi(Z_v(y))\), where \(Z_v\) is that future
Langevin flow.  On a good path from \(x\), cut this test off on the
inverse-polynomial ball just described.  Its gradient there is bounded by
\(e^{-v/2}\|\nabla\phi(Z_v(y))\|\).  Crossing and cutoff-annulus errors
are negligible by short-time Gaussian tails after making \(\Delta t\)
small enough.

The good-path event is determined from the future path started at
\(x\), and the cutoff test has polynomial bounds on this ball.
Bad paths, averaged over stationary \(x\), cost only a polynomial
times their probability. On good paths, drop the event indicator
from the resulting nonnegative gradient integral. Then
\(x\sim P\), \(y\sim K_{\Delta t}(x,\cdot)\), and independent
future noise make \(Z_v(y)\) stationary. Each variance term is
therefore at most
\[
 (2+o(1))\Delta t\,e^{-v}
       \int\|\nabla\phi\|^2dP
\]
plus an arbitrarily small inverse-polynomial error.  Summing over the polynomial number of time steps preserves the
negligible error. The gradient contributions sum to at most
\((2+o(1))\int\|\nabla\phi\|^2dP\).

It remains to bound \(\operatorname{Var}_P K_{T'}\phi\).
Let \(b(x)\) be the conditional probability that the future path
from \(x\) fails the preceding tube test. Its \(P\) mean is at
most \(\operatorname{poly}(n)e^{-\beta k}+e^{-n^2}\). Markov at
\(e^{-\beta k/3}\), using \(\beta k\gg\log n\), leaves a set of
\(P\)-mass at most \(e^{-\beta k/2}\). Off that set, averaging the
common-noise coupling gives local differences bounded by
\(e^{-T'/2}\operatorname{Lip}(\phi)\) times distance, plus an
arbitrarily small inverse-polynomial error.  Transfer the exceptional
starting set to \(G\) using \eqref{prof:eq-likelihood-lower-tail}.
For independent \(g_0,g_1\sim G\), use a polynomial mesh on
\(\cos\theta\,g_0+\sin\theta\,g_1\), \(0\le\theta\le\pi/2\).
Every vertex has law \(G\), and Gaussian norm tails put adjacent
vertices within the local ball. A union bound over the vertices
leaves Gaussian exceptional probability
\(\operatorname{poly}(n)e^{-c\beta k}\).

The local estimates telescope along this polygon. To reweight its
endpoints to two independent \(P\)-points, Cauchy--Schwarz and
\eqref{prof:eq20} bound the exceptional probability by
\[
 e^{\alpha k}
       \bigl(\operatorname{poly}(n)e^{-c\beta k}\bigr)^{1/2}.
\]
It is smaller than any required inverse polynomial when
\(\alpha/\beta\) is sufficiently small and \(\beta k/\log n\)
is sufficiently large. Taking \(C_D\) large gives
the desired residual variance bound.  Absorb the remaining errors in
the slack between \(2+o(1)\) and \(3\).
\end{proof}

\subsection{Backward amplification to logarithmic observation scale}
\label{prof:sec-amplification}

We use two positive-time inputs: for every compact
\([s_1,T]\subset(0,\infty)\), posterior covariance is bounded by a
constant except at integrated planted probability \(e^{-c(s_1,T)n}\);
and at a sufficiently large fixed \(T\), the posterior unscaled gap is
bounded below with the same type of exception.  The critical endpoint
verification of these inputs is given in the positive-time argument.
Their use here is only with the rotation-invariant restriction
\(\mathcal G\), and only on fixed positive intervals.

\begin{proposition}[Amplified covariance estimate]
\label{prof:prop-amplified-covariance}
For every sufficiently small fixed \(\epsilon>0\), there are
\(s_0,a>0\) and \(K_{\log}<\infty\) such that, at every deterministic
time in the indicated range,
\begin{equation}
 \|\Sigma_s\|\le\frac{1+\epsilon}{s},
 \qquad 0<s\le s_0,\qquad ns^{3/2}\ge K_{\log}\log n,
 \label{prof:eq22}
\end{equation}
outside an integrated planted exceptional set, intersected with
\(\mathcal G\), of mass at most \(e^{-ans^{3/2}}\).
For fixed \(1/2<\gamma<\gamma'<1\), the same exceptional set may be
used for
\begin{equation}
 \begin{split}
 \|A_0\Sigma_sA_0\|&\le r^{-\gamma},\qquad A_0=2I-J,\\
 v_s&:=\sup_{\substack{H^T=-H\\ \|H\|_{\mathrm{HS}}\le1}}
 \operatorname{Var}_{\mu_{h_s}}
 \bigl((X-m_s^{\mathrm{post}})^T[H,J](X-m_s^{\mathrm{post}})\bigr)
       \le C_v s^{-1-\gamma'/2}.
 \end{split}
 \label{prof:eq23}
\end{equation}
\end{proposition}

\begin{proof}
Each backward step first uses later observations to obtain coarse
covariance and moment bounds. These control derivatives along rotation
orbits, where the static majority estimates reset covariance to its sharp
threshold. The initial positive-time band uses the same reset; its coarse
bounds come from the separate fixed-positive-time and terminal inputs.
We carry all three estimates together. First fix
\(0<\iota<\gamma'/2\), with the target covariance slack
\(\epsilon\) small enough for the variance calculation below.
That calculation gives a constant \(C_\iota'\) independent of
\(L_0\) and of the relative step size. Choose \(L_0\) large enough
that \(C_\iota'L_0^{\iota-\gamma'/2}<1/4\), and then choose
\(s_0\) small enough for the static medians and
Lemma~\ref{prof:lem-compact-bridge}.

The quadratic-variance constant
\(C_v\) and the auxiliary fourth-moment constants below may depend on
these choices, but will not depend on the eventual relative step size.
The initial band is \([s_0/L_0,s_0]\). Let \(F_b^+\) denote the
positive-time covariance failure, including the terminal gap failure
when \(b=T\), and define on this band
\[
 Q_s^+=\mathbf1_{F_s^+}+\Pr_s(F_T^+)
       +\int_s^T(1+b^{-C})\Pr_s(F_b^+)\,db.
\]
Here and below \(\Pr_s\) conditions on \((J,h_s)\), and the fixed
\(C\) dominates the weights in the moment transfers. Tonelli and
Markov on the fixed interval \([s_0/L_0,T]\) give a fixed
\(c_{\rm init}>0\) such that \(Q_s^+\le e^{-c_{\rm init}n}\)
outside integrated mass \(e^{-c_{\rm init}n}\). On this one event,
the bounded-covariance martingale transfer to \(T\) and its gap
bound give all the initial quadratic and fourth-moment bounds proved
below. No conditional variance lemma is needed for this initial
transfer. Orientation concentration will reset the covariance
thresholds on the same band.

Let \(c_+>0\) be a positive-time failure rate on \([s_0,T]\),
and let \(a_{\rm stat}>0\) be smaller than the fixed root,
orbit, diagnostic, and spherical-width exclusion rates at the
accuracies used below. The compact bridge has its endpoint-independent
rate \(a_*\). Choose \(a>0\) once, with
\[
 8a<\min\{a_{\rm stat},a_*,
          c_{\rm init}/s_0^{3/2},c_+/s_0^{3/2},1\},
\]
and small enough for the final orbit-gradient condition whenever
its concentration exponent is at most \(2a\). That condition uses
only the moment constants and the already fixed \(s_0,L_0\).

\paragraph{A conditional event with room in its probability exponent.}
For \(s<s_0/L_0\), put \(u=(1+d)s\) and \(t=L_0s\), and suppose the
induction estimates hold at all required times at least \(u\).
The active \(d\) will be one of two fixed values specified below;
the rate \(a\) will not be decreased after those choices.
Let \(F_b\) be the union of the three inductive failures at
\(b\le s_0\), and the positive-time covariance failure at
\(s_0<b\le T\). Include the terminal gap failure in \(F_T\).
Define
\[
 Q_s=\Pr_s(F_u)+\Pr_s(F_t)+\Pr_s(F_T)
       +\int_u^T(1+b^{-C})\Pr_s(F_b)\,db.
\]
Choose the fixed exponent \(C\) large enough to dominate the
polynomial weights in the variance and fourth-moment estimates
below. All these weights are polynomial in \(n\) on the stated
time range. Tonelli's theorem and the deterministic-time induction
give
\[
 \E_{\rm pl}[\mathbf1_{\mathcal G}Q_s]
 \le \operatorname{poly}(n)e^{-ak(u)}
       +\operatorname{poly}(n)e^{-c_+n}.
\]
Here \(\mathcal G\) is measurable at time \(s\); the equality
behind Tonelli therefore retains its indicator. Above \(s_0\) we
used the positive-time inputs on the fixed interval \([s_0,T]\).

Set \(\beta=c'da\), with a sufficiently small absolute \(c'>0\).
Markov at \(Q_s=e^{-\beta k}\) gives, before polynomial costs, the
exponent
\[
 \bigl(a(1+d)^{3/2}-\beta\bigr)k.
\]
For small fixed \(d\) this strictly exceeds \(ak\).
The other Markov term is
\(\operatorname{poly}(n)e^{-c_+n+\beta k}\); the prior choice of
\(a\), and \(k\le ns_0^{3/2}\), also place it strictly above
the required rate. We next intersect this test with the likelihood
and diagnostic events needed for the derivative bounds.

For sufficiently small \(d\), the covariance hypothesis in
\eqref{prof:eq20} follows from \(Q_s\le e^{-\beta k}\).
The likelihood hypothesis
follows from the exact Gaussian identity
\[
 \int e^{2L_P}dG
 =\mathbb E_s\exp\{\tau(X^1-m_{\mathrm{sh}})^T
                              (X^2-m_{\mathrm{sh}})\}.
\]
The initial band is already fixed, so the back-step must be chosen
without shrinking that band again. Two likelihood estimates provide
this separation of choices: the spherical log-moment estimate near zero,
and the compact overlap bridge on the remaining interval.

There are two parameter ranges. At sufficiently small \(r\), use the
spherical shift \(m_{\mathrm{sh}}=m_o\). Choose
\(d_{\rm near}>0\) so that \(Cd_{\rm near}^2\) is a
sufficiently small multiple of \(d_{\rm near}a\). Apply the
off-diagonal replica log-moment bound~\eqref{prof:eq5} to this
fixed two-replica test, with requested failure rate \(4a\).
Given \(\varepsilon_{\rm like}>0\), it gives
\[
 \log\int e^{2L_P}dG
 \le Cd_{\rm near}^2k+\varepsilon_{\rm like}k+C_2\log n
\]
on a sufficiently small upper range. Choose
\(\varepsilon_{\rm like}\ll d_{\rm near}a\), then its upper
endpoint \(r_c>0\), and finally take \(K_{\log}\) large enough
that \(C_2\log n\ll d_{\rm near}ak\). Shrinking \(r_c\)
alone does not absorb this logarithmic term.
Only this selected test is needed, not simultaneous control for
all prospective step sizes. The width exclusion has the larger
rate fixed above. On the remaining compact
interval use \(m_{\mathrm{sh}}=m\) and
Lemma~\ref{prof:lem-compact-bridge}, with \(\chi\ll a\).

Choose a second step \(d_{\rm comp}>0\) on this compact interval.
The good-pair contribution to the second moment is at most
\(e^{d_{\rm comp}\chi k}\), while the exceptional-pair
contribution is at most
\[
 e^{-cn+4d_{\rm comp}s n},
\]
using \(\|X-m\|\le2\sqrt n\). Choose \(d_{\rm comp}\) after
the compact-tail constant \(c\) so that this term is negligible.
Then \(\alpha\le 2d_{\rm comp}\chi\) for large \(n\), while
\(\beta=c'd_{\rm comp}a\). In either range,
\(\alpha/\beta\) in \eqref{prof:eq20} is as small as required.
Thus no choice of \(s_0\) depends on \(a\) or on these small step sizes.
Choose \(K_{\log}\) last to absorb the polynomial costs of
Lemma~\ref{prof:lem-conditional-variance}.

For each active step let \(\mathcal L_s\) be its likelihood event
just constructed, and let \(\mathcal D_s\) retain the diagnostics
and root or spherical-width events used in its orbit comparison.
The full event used from now on is
\[
 \mathcal C_s=\{Q_s\le e^{-\beta k}\}
                    \cap\mathcal L_s\cap\mathcal D_s.
\]
The first two factors are tests of \((J,h_s)\), and are the only
ones used to apply the conditional variance lemma under \(\Pr_s\).
The diagnostic factor may also use the plant and root-orbit data;
it restricts the subsequent orbit integration without further
conditioning the posterior spin law. Put
\[
 \Delta=\min_{d\in\{d_{\rm near},d_{\rm comp}\}}
       \{a((1+d)^{3/2}-1)-c'da\}>0.
\]
The preceding estimates give
\[
 \Pr_{\rm pl}(\mathcal G\cap\mathcal C_s^c)
 \le\operatorname{poly}(n)e^{-(a+\Delta)k}
       +\operatorname{poly}(n)e^{-4ak}.
\]
All assertions remain integrated with \(\mathcal G\) imposed;
none is a bound conditional on every spectrum. On the initial band
define \(\mathcal C_s\) instead as
\(\{Q_s^+\le e^{-c_{\rm init}n}\}\cap\mathcal D_s\),
whose excluded mass has a fixed rate larger than \(4a\).

\paragraph{Coarse covariance and a longer variance transfer.}
On \(\mathcal C_s\), total covariance gives
\[
 \Sigma_s=\mathbb E_s\Sigma_u+
                  \operatorname{Cov}_s(m_u^{\mathrm{post}}).
\]
Apply \eqref{prof:eq21} to every linear projection of
\(m_u^{\mathrm{post}}\).  Its derivative in \(g\) is
\(\sqrt\tau\,\Sigma_u\), so the second term is at most
\(3\tau\mathbb E_s\Sigma_u^2\), up to an arbitrarily small
inverse-polynomial operator error.  Apply the same argument after
multiplication by \(A_0\).  This gives both covariance bounds with
fixed coarse factors.

For any smooth polynomially bounded test of the later field,
\begin{equation}
 \operatorname{Var}_s\bigl(\phi(h_t)\bigr)
 \le C_\iota tL_0^{1+\iota}
             \mathbb E_s\|\nabla\phi(h_t)\|^2+n^{-D}.
 \label{prof:eq24}
\end{equation}
Here \(\iota>0\) is arbitrarily small if \(\epsilon\) is chosen
sufficiently small, and \(C_\iota\) is independent of \(L_0,d\).
To prove this, the posterior-expectation martingale on \([u,t]\) has
diffusion coefficient
\[
 \nabla\mathbb E_b\phi(h_t)
  =\mathbb E_b\nabla\phi(h_t)
       +\operatorname{Cov}_b(X,\phi(h_t)).
\]
This follows directly from the Gaussian likelihood of the future
observation increment given \(X\).

The observations have unit
diffusion and posterior-mean drift in their own filtration.  If
\(R_b=\mathbb E_s\operatorname{Var}_b(\phi(h_t))\), their martingale
variance identity and Cauchy--Schwarz yield
\[
 -R_b'\le \frac{1+\iota}{b}R_b
       +C_\iota\mathbb E_s\|\nabla\phi(h_t)\|^2
       +\operatorname{poly}(n)\Pr_s(\text{failure at }b).
\]
In the covariance term use
\(\|\operatorname{Cov}_b(X,Y)\|^2\le
\|\Sigma_b\|\operatorname{Var}_bY\), and absorb the mixed term
by the slack in \(\iota\).

Integrating backwards from \(R_t=0\) gives explicitly
\[
 R_u\le C_\iota t(t/u)^{1+\iota}
       \E_s\|\nabla\phi(h_t)\|^2
   +\int_u^t (b/u)^{1+\iota}
       \operatorname{poly}(n)\Pr_s(F_b)\,db.
\]
The last term is negligible on the conditional event. For the first interval
\([s,u]\), apply \eqref{prof:eq21} to \(\mathbb E_u\phi(h_t)\)
and use the same gradient estimate. It adds at most a fixed multiple
of \(R_u+\tau\E_s\|\nabla\phi(h_t)\|^2\), since \(\tau/u\le d\).
As \(t/u\le L_0\), this proves \eqref{prof:eq24} with a constant
independent of \(L_0,d\).
When covariance is bounded on a fixed positive band beginning at \(s\),
this same martingale argument works directly from \(s\) up to \(T\),
with fixed constants in place of the displayed powers.

\paragraph{Quadratic and fourth moments.}
Let \(M=[H,J]=-[H,A_0]\), with \(\|H\|_{\mathrm{HS}}\le1\), and
put \(D_t=m_t^{\mathrm{post}}-m_s^{\mathrm{post}}\).  Total covariance
gives \(\mathbb E_sD_tD_t^T\preceq\Sigma_s\).  At time \(t\),
the centered quadratic from time \(s\) splits into a centered quadratic
there and the linear term
\(2D_t^TM(X-m_t^{\mathrm{post}})\), plus its conditional mean.
The averaged variance of this linear term is
\(O_{L_0}(s^{-1-\gamma/2})\). Indeed expand
\(M=-HA_0+A_0H\). On successful descendants the two terms obey
\[
\begin{split}
 \E_s[\mathbf1_{\rm good}D_t^THA_0\Sigma_tA_0H^TD_t]
 &\le Ct^{-\gamma/2}\operatorname{Tr}(H^T\Sigma_sH)
 \le Ct^{-\gamma/2}s^{-1},\\
 \E_s[\mathbf1_{\rm good}D_t^TA_0H\Sigma_tH^TA_0D_t]
 &\le Ct^{-1}\operatorname{Tr}(H^TA_0\Sigma_sA_0H)
 \le Ct^{-1}s^{-\gamma/2}.
\end{split}
\]
The norm bound at \(t\) is applied before averaging, and
\(\E_sD_tD_t^T\preceq\Sigma_s\) is then used in the trace.
No independence between \(D_t\) and \(\Sigma_t\) is needed.
Polynomially bounded errors on failed descendants are negligible.

The conditional mean of the whole quadratic is
\(\operatorname{Tr}(M\Sigma_t)+D_t^TMD_t\).  As a function of the
field at \(t\), its gradient is its conditional covariance, as a spin
test, with \(X\).  On successful descendants its squared norm is thus
at most \((1+\epsilon)/t\) times its conditional variance.  Elsewhere
use the polynomial deterministic bound.  Total variance,
\eqref{prof:eq24}, and splitting sums with a factor two give
\[
 v_s\le C_\iota' L_0^{1+\iota}C_vt^{-1-\gamma'/2}
           +C_{L_0}s^{-1-\gamma/2}+o(1).
\]
The coefficient of the first term relative to
\(C_vs^{-1-\gamma'/2}\) is
\(C_\iota'L_0^{\iota-\gamma'/2}<1/4\) by our earlier choice.
Choose \(C_v\) large. The strict
inequality \(\gamma<\gamma'\) absorbs the second term.  This closes
the quadratic-variance induction.

On the initial band, apply the bounded-covariance version of the transfer
to \(T\).  At \(T\), the gap input bounds the centered quadratic
variance because its single-site heat-bath energy is at most a constant
times
\(\operatorname{Tr}(M\Sigma_TM)+\|M\|_{\mathrm{HS}}^2\),
as follows by expanding a single spin flip.  Coarse covariance on the
initial band comes directly from the positive-time input.

We also obtain uniformly in unit \(e\)
\[
 \mathbb E_s|e^T(X-m_s^{\mathrm{post}})|^4\le Cr^{-4},\qquad
 \mathbb E_s|e^TA_0(X-m_s^{\mathrm{post}})|^4\le Cr^{-2\gamma}.
\]
Integrate the posterior-mean martingale from \(u\) to \(T\), whose
diffusion coefficient is \(\Sigma_b\).  The respective directional
quadratic variations have second moments bounded by the squares of
the integrals of \(Cb^{-2}\) and \(Cb^{-1-\gamma/2}\) on small
scales; use bounded covariance later.  These estimates follow from
\(\|\Sigma_b\|^2\) and
\(\|\Sigma_bA_0\|^2\le
\|\Sigma_b\|\|A_0\Sigma_bA_0\|\), with deterministic polynomial
bounds on exceptions.

Martingale fourth-moment inequalities apply.
Terminal fourth moments follow by applying the gap first to a linear
test and then to its centered square.  Finally apply
\eqref{prof:eq21} to a centered linear projection of
\(m_u^{\mathrm{post}}\) and to its square; their gradients involve
\(\Sigma_u\).  This treats
\(m_u^{\mathrm{post}}-m_s^{\mathrm{post}}\) without constants
depending on \(d\).  On the initial band start the mean martingale
directly at \(s\).

\paragraph{Restoring the sharp covariance thresholds.}
On each root orbit, monitor the two covariance norms and the two mean
deviations in the static comparison, in their own normalized units.
Use tolerance \(\epsilon/2\) for the sharp covariance median and
\(\epsilon\) for its failure threshold, with fixed slack also for
the other quantities.  In particular the mean thresholds are
\(r^{-1-\eta}\) and \(r^{-\gamma/2}\), with an arbitrarily small
\(\eta>0\).  Let a Lipschitz ramp be zero below the median thresholds
and one when any failure threshold is reached.

On the conditional
events already constructed, its gradient in the transition region is
bounded by
\[
 D_{\mathrm{grad}}=Cr^{-1-\gamma''/2},\qquad
                     \gamma'<\gamma''<1.
\]
To check this, a root-preserving rotation has score \(z^TMz/2\), with
\(m\) fixed.  Its centered quadratic variance is bounded by
\eqref{prof:eq23}.  The additional linear term about
\(m_s^{\mathrm{post}}\) has variance at most
\(Cr^{-2-\gamma-2\eta}\), by the coarse covariance bounds and the
transition thresholds.

Choose \(\eta\) sufficiently small.
Cauchy--Schwarz with the fourth and second moments just proved then
controls derivatives of directional centered squares and means.
The derivative of \(A_0\) itself costs at most
\(Cr^{-1-\eta}\) in these units.  Outside these conditional events
the ramp has a polynomial global gradient bound.  This calculation
does not differentiate the root.

When \(K_{\log}\log n\le k<\log^{10}n\), use instead unconstrained
orbits that co-rotate the field and matrix, centered at \(m_o\).
The polylogarithmic edge comparison and the spherical shell estimate
give the required medians, after excluding the small-rate event on
which the spherical width squared differs from \(s\) by more than a
fixed relative tolerance \(\epsilon'\ll\epsilon\). Here is the
conversion from raw moments. In this band
\(r=n^{-1/3}\log^{O(1)}n\), so the unweighted error
\(n^{1/2+o(1)}\) in \Cref{prof:prop:polylog} is \(o(s^{-1})\).
On the width event, \eqref{prof:eq2} bounds the spherical raw
moment about \(m_o\) by
\((1+O(\epsilon')+o(1))s^{-1}I\). The cube covariance is bounded
by its raw moment about \(m_o\), giving the sharp covariance
median. The weighted raw moment is at most \(n^{1/6+o(1)}\),
which is \(o(r^{-\gamma})\) because \(\gamma/3>1/6\).
Jensen bounds the squares of the ordinary and weighted mean
deviations by their respective raw moment operators. These give
the two looser mean medians \(r^{-1-\eta}\) and
\(r^{-\gamma/2}\). All four statements concern a Haar majority
on the orbit that co-rotates this field.

The extra linear
score coefficient is
\[
 Mm_o+Hh=(\lambda-J)Hm_o.
\]
Since \(\|m_o\|=O(\sqrt k/r)\), its standard deviation has an
additional multiplier at most \(C\sqrt k\) in the preceding bounds.
Differentiating the center costs the same allowance.  Thus
\(D_{\mathrm{grad}}\) may be enlarged by a fixed power of \(\log n\)
in this range.  The unconstrained partition-ratio comparison~\eqref{prof:eq4} transfers
the orbit weights.

The conditional derivative bound has exceptions whose masses depend
on the orbit. We need a concentration estimate that is linear in those
masses, so that it can be averaged over the orbit parameters.  A Haar ramp zero on a
fixed majority, with gradient at most \(D_{\mathrm{grad}}\) off a set
of mass \(p\), and polynomially bounded gradient everywhere, has
plateau mass at most
\[
 e^{-bk}+\operatorname{poly}(n)p,
\]
provided \(D_{\mathrm{grad}}\sqrt{bk/n}\) is sufficiently small.
Indeed split the upper ramp levels at tail probabilities successively
halving from \(1/2\) to \(e^{-bk}\).

If the plateau exceeds both
this last probability and a large polynomial times \(p\), truncate
and rescale the ramp across one such quantile interval of width \(w\)
and probability level \(p'\).  The entropy of its square is at least
\(cp'\log(1/p')\), whereas its gradient-square integral is at most
\[
 w^{-2}\bigl(2D_{\mathrm{grad}}^2p'
                         +\operatorname{poly}(n)p\bigr).
\]
The rotation-group heat log-Sobolev inequality gives
\(w\le CD_{\mathrm{grad}}/\sqrt{n\log(1/p')}\).

Strict supports and inclusive plateaux handle atoms or zero-width
intervals.  Summing these widths bounds the full rise by
\(CD_{\mathrm{grad}}\sqrt{bk/n}\), a contradiction.
In particular the orbit-dependent exceptional masses enter linearly
and can be averaged without taking a fractional power.

Take \(b=a+\Delta/2<2a\). The gradient condition uses only
\(L_0,s_0,C_v\), the moment constants and \(\gamma''<1\), and
holds by the earlier choice of \(a\). It also
holds throughout the polylogarithmic band despite the extra logarithms.
On regular root anchors, the determinant-weight comparison at the
concentration radius costs
\[
 O\bigl(\sqrt k+k/(nq^2)\bigr)=o(k),
\]
so it does not require shrinking \(r\) after the rate has been chosen.

We display the two measure changes. Let \(p(\omega)\) be the
orbit mass of \(\mathcal C_s^c\), including its likelihood and
diagnostic exclusions, and let \(h(\omega)\) be
its ramp plateau mass. The quantile estimate and
\(\int w_0\,d\pi\le n^C\) give
\[
 \int w_0h\,d\pi
 \le e^{-bk}\int w_0\,d\pi
       +\operatorname{poly}(n)\int w_0p\,d\pi.
\]
The first comparison in \Cref{prof:orbits} bounds
\(\int w_0p\,d\pi\) by the counting derivative-failure mass
times \(e^{\xi k}\), plus larger-rate exclusions. Returning the
last display to the counting law by the reverse comparison costs
a second \(e^{\xi k}\). Thus the derivative term retains exponent
at least \(a+\Delta-2\xi\), and the direct plateau term retains
exponent at least \(b-\xi\), before polynomial factors. Choose
\(0<\xi<\Delta/8\). The excluded
diagnostic and base-orbit terms have already been assigned larger
rates. Finally the degree identity charges excess root multiplicity
to the excluded negative-Jacobian roots when returning from counts
to planted fields.
For unconstrained orbits choose partition-ratio accuracies relative to
the small-range step size, then enlarge \(K_{\log}\).  Every reset
cost and discarded event therefore fits within \(e^{-ak}\).
The same reset works on the initial band with its fixed constants
and the larger initial-event rate already reserved; \(a\) is unchanged.

For clarity, the order of choices is: the target accuracies,
\(\iota<\gamma'/2\), and then large \(L_0\); small \(s_0\) for the static medians, the fixed outer cuts,
and the compact bridge; then \(C_v\) and all moment constants,
including those on the initial band; then \(a\); then the two
positive relative step sizes \(d_{\rm near},d_{\rm comp}\); and
finally \(K_{\log}\). Put
\(d_{\min}=\min(d_{\rm near},d_{\rm comp})>0\).
On successive backward bands with ratio \(1+d_{\min}\), every
active \(u=(1+d(s))s\) lies in an earlier band, as does \(t\).
Induction on these bands proves \eqref{prof:eq22}--\eqref{prof:eq23}
for each deterministic \(s\). Tonelli supplies the time integrals
used above from these pointwise statements; no union over
uncountably many observation times is taken. No positive-time input
is required uniformly down to an endpoint depending on \(a\).
\end{proof}


\subsection{Covariance below the logarithmic threshold}
\label{prof:sec-bottom-scale}

The backward induction stops at logarithmic \(k\). Below that threshold
we prove a separate, coarser covariance bound with an arbitrarily large
polynomial probability exponent. The required probability cannot be
obtained by merely repeating the second-moment comparison.  The Haar Gram comparison and the mixed
cube--sphere tensor cancellation are closely related to
\citet[Section~3.3]{DuHuang2026CriticalFluctuations} and
\citet[Lemma~3.5]{DuHuang2026CriticalOverlap}.  Here the full
finite-replica expansion keeps track of every normalization factor to
obtain the required probability exponent along the observation fields.

\begin{proposition}[Bottom-scale covariance]
\label{prof:prop-bottom-covariance}
For every fixed \(B>0\) and \(K_{\log}<\infty\), at each deterministic
time satisfying \(ns^{3/2}\le K_{\log}\log n\),
\begin{equation}
 \|\Sigma_s\|\le Cn^{2/3}\ell^C,
 \qquad \ell=\log\log(n+1000),
 \label{prof:eq25}
\end{equation}
except on an integrated planted event, intersected with
\(\mathcal G\), of mass at most \(n^{-B}\).
\end{proposition}

\begin{proof}
Set \(L=n^{2/3}\).  The spherical raw second moment about \(m_o\)
already satisfies the bound in \eqref{prof:eq25}.  The hypotheses of the
polylogarithmic edge comparison, namely
\[
 \lambda-\lambda^0\le n^{-2/3}\log^C n,
 \qquad \|h\|\le\sqrt n\,r\log n,
\]
may be imposed with arbitrarily small inverse-polynomial failure by
testing the spherical observation at a larger polylogarithmic edge
width.  We prove a higher-moment version of that comparison.  Only
unweighted raw second moments are needed here.

On an unconstrained Haar orbit divide every integral by \(Z_o(h)\).
Let \(\Delta\) be either the signed cube-minus-sphere mass, or the signed
matrix integral of \((x-m_o)(x-m_o)^T/L\), with this common denominator.
For a fixed even integer \(p\), expand \(\|\Delta\|_{\mathrm{HS}}^p\)
(the mass case means \(|\Delta|^p\)).  A subset \(B'\subseteq[p]\)
specifies which of the \(p\) draws have sign priors; the other draws have
spherical priors.  The common kernel is respectively \(1\), or
\[
 \prod_{a=1}^{p/2}
 \left(\frac{(x_{2a-1}-m_o)^T(x_{2a}-m_o)}L\right)^2.
\]
The sign of that term is \((-1)^{p-|B'|}\).

\paragraph{Truncating the moment before annealing.}
In each term first restrict the internal sign overlaps
\(Q_{ij}=x_i^Tx_j/n\), \(i,j\in B'\), to
\(|Q_{ij}|\le\eta_p\), where \(\eta_p>0\) is fixed and small.
This avoids assuming an unrestricted high moment of the partition
function.  Outside exponentially exceptional orientations the truncated
alternating sum differs negligibly from the desired power: the
unnormalized large-overlap sign-pair integral is exponentially small by
the global pair bound, while the total normalized cube mass is at most
\(e^{c_{\mathrm{tiny}}n}\) by its first moment.  Choose
\(c_{\mathrm{tiny}}\) sufficiently small after \(p,\eta_p\).
The polynomial kernel costs can also be absorbed.

To discard the exceptional orientations inside an alternating sum,
we need a bound on each term separately. Let \(T_{B'}(O)\ge0\)
be the normalized, individually truncated term at orientation \(O\),
including its mass or product-of-squares kernel. We claim
\[
 \E_{\rm Haar}T_{B'}^2\le e^{o(n)}.
\]
To verify this, bound the polynomial insertions crudely
and integrate out the spherical priors.  Two groups of sign vectors
remain, with nearly identity internal Grams.  If \(R\) is their matrix
of cross overlaps, the empirical-site-type entropy cost is at least
\[
 \frac{\|R\|_{\mathrm{HS}}^2}{2(1+C_p\eta_p)}
                      +c_p\|R\|_{\mathrm{HS}}^4.
\]
Indeed let \(E\) be the first-group sign vector in the empirical
site-type law, and let \(b_j(E)\) be the conditional bias of the
\(j\)th sign in the second group. Then
\[
 R_{ij}=\E[E_i b_j(E)],\qquad
 \sum_j\E b_j(E)^2
       \ge \frac{\|R\|_{\rm HS}^2}{1+C_p\eta_p},
\]
by projection onto the first-group coordinates and their Gram bound.
Conditional relative entropy dominates the sum of marginal
conditional entropies. The binary entropy inequality
\(I(b)\ge b^2/2+c b^4\), followed by Jensen and the fixed number
of coordinates, gives the displayed quadratic and quartic costs.
On the other hand, diagonalizing the
full Gram and applying the semicircle angular bound successively shows
that the angular energy gain above the independent baseline is at most
\(n/4\) times the squared Hilbert--Schmidt norm of the full off-diagonal
Gram, plus \(o(n)\). The cross block \(R\) occurs twice in that
norm, so the angular gain is bounded above by
\[
 \frac n2\|R\|_{\rm HS}^2+O_p(n\eta_p^2)+o(n).
\]
The leading quadratic halves cancel against the entropy cost.
The remaining positive terms are
\(C_p n\eta_p\|R\|_{\rm HS}^2+O_p(n\eta_p^2)\), whereas the
quartic term is \(-c_p n\|R\|_{\rm HS}^4\). Thus cross overlaps outside any sufficiently small
fixed neighborhood of zero have a strict negative rate after choosing
\(\eta_p\) sufficiently small.  Within that neighborhood the small-Gram
comparison and the spherical length-deficit tails give \(e^{o(n)}\).
This proves the second-moment bound. Cauchy--Schwarz applied to
each \(T_{B'}\) removes an exponentially exceptional orientation
set at cost \(e^{o(n)}e^{-cn/2}\). It may therefore be removed
before the signed sum, even though that sum need not be nonnegative
on the exceptional set.

The same estimates restrict each term's Haar integral, to arbitrarily
high inverse-polynomial accuracy, to
\begin{equation}
 |Q_{ij}|\le (L/n)\log^{C_p}n\quad(i\ne j).
 \label{prof:eq-bottom-window}
\end{equation}
For internal sign overlaps, use the small-Gram comparison until this
window and choose its logarithmic power large enough that the
length-deficit decay absorbs all losses.  Once those overlaps are in the
window, all overlaps involving spherical priors have the same tails:
conditional on the internal sign Gram, the rotated prior is exactly the
spherical prior conditioned on that Gram.  Polynomially fine boxes have
controlled angular variation on these nonsingular internal Grams.
The central cutoffs may be smooth.  The large-overlap truncation has
therefore disappeared before the alternating cancellation below.

\paragraph{The finite-label expansion.}
The overlap window is now common to all terms. The only remaining
prior ratio is the internal sign Gram mass, in lattice density units,
divided by the spherical internal Gram density on the label set \(B'\).
We expand this ratio with every normalization factor included. Each
coefficient must have an explicit finite set of replica labels: the
alternating sum cancels exactly those terms that omit a label.

For \(B'\subseteq[p]\), put \(v=|B'|\), \(d_v=\binom v2\), and let
\(Q_{B'}\) be the Gram with diagonal one.  Zero-dimensional densities
and empty products are one.  Let \(\rho_{B'}\) be the preceding density
ratio.  To every fixed expansion order it has the form
\begin{equation}
 \rho_{B'}(Q,n)=
 \exp\left\{n\sum_{b\ge4}F_{B',b}(Q)\right\}
 \sum_{b,j\ge0}n^{-j/2}A_{B',b,j}(Q),
 \qquad A_{B',0,0}=1.
 \label{prof:eq-label-expansion}
\end{equation}
Here the displayed coefficients are homogeneous of degree \(b\) in the
off-diagonal overlaps.

Every degree-\(b\) rate coefficient is a sum of
terms involving at most \(6b\) replica labels; every amplitude coefficient
of degree \(b\) and order \(n^{-j/2}\) is such a sum involving at most
\(6(b+j)\) labels.  A term involving a label set \(S\) means a universal
summand indexed by labels in \(S\), including all internally summed
indices, whose value is independent of whether further labels belong to
\(B'\).  This qualification includes the normalized central density
prefactors, rather than merely the visible overlap monomial.  On
\eqref{prof:eq-bottom-window}, the expansion and its exponential can be
truncated with uniform relative error \(O(n^{-A})\), for any prescribed
fixed \(A\).

We prove these assertions, starting with the normalizations.  For one
sign site let
\(Y=(\varepsilon_a\varepsilon_b)_{a<b,\ a,b\in B'}\).
If \(\delta(S)\) denotes the cut vector of a subset of labels, then
\(Y=\mathbf1-2\delta(S)\).  Its difference lattice is
\(\Lambda_v=\operatorname{span}_{\mathbb Z}\{Y-\mathbf1\}\).
Distinguish one label, denoted \(i_\bullet\) in this lattice calculation.
A basis consists of
\[
 \{2\delta(\{i\}):i\ne i_\bullet\}
 \ \cup\
 \{4e_{ij}:i<j,\ i,j\ne i_\bullet\}.
\]
Indeed
\[
 \delta(S)=\sum_{i\in S}\delta(\{i\})
       -2\sum_{\{i,j\}\subseteq S}e_{ij},\qquad
 \delta(\{i\})+\delta(\{j\})-\delta(\{i,j\})=2e_{ij}.
\]
Ordering the edges incident to \(i_\bullet\) first makes the basis block
triangular, so its covolume is
\[
 c_v=2^{v-1}4^{\binom{v-1}2}=2^{(v-1)^2}\quad(v\ge1),
 \qquad c_0=1.
\]
The sample overlap belongs to \(\mathbf1+n^{-1}\Lambda_v\), whose
cell volume is \(c_vn^{-d_v}\).

Thus the sign density in lattice units is
\[
 f^S_{n,B'}(Q)=\frac{n^{d_v}}{c_v}
    \Pr\left\{n^{-1}\sum_{i=1}^nY_i=(Q_{ab})_{a<b}\right\}.
\]
The reciprocal lattice, using the \(2\pi\) convention, is
\[
 \Lambda_v^\vee=\frac\pi2
 \left\{z\in\mathbb Z^{d_v}:
           \sum_{b\ne a}z_{ab}\equiv0\pmod2
                           \text{ for every }a\right\}.
\]
Equality in the triangle inequality shows that the characteristic
function of \(Y\) has modulus one precisely on this reciprocal lattice.
The same is true under every sufficiently small real exponential tilt,
whose support is unchanged.  There is consequently a single central
saddle on the reciprocal torus, and uniform contraction away from it.
Boundary representatives of that one saddle are identified on the torus.

For the spherical density, use independent Gaussian coordinates and
the one-site statistics
\[
 \left((g_ag_b)_{a<b},\ ((g_a^2-1)/\sqrt2)_{a\in B'}\right).
\]
Their covariance is the identity in dimension \(d_v+v\); the sign
statistics also have identity covariance in dimension \(d_v\).
Denote the two log transforms by \(\Psi_G,\Psi_S\), and their Legendre
transforms by \(I_G,I_S\).  At target \((Q_{ab},0)\),
\(I_G=-\tfrac12\log\det Q_{B'}\).  If \(f^G\) is the full Gaussian
statistic density, conditioning on the zero normalized-length targets
gives exactly
\[
 f^O_{n,B'}(Q)=f^G_{n,B'}(Q,0)/f^{\mathrm{len}}_{n,B'}(0).
\]
There is no extra Jacobian because both densities use the same normalized
length coordinates.

With \(z_n=n/2\), their exact length-only factor is
\[
 f^{\mathrm{len}}_{n,B'}(0)
    =\left(\frac n{2\pi}\right)^{v/2}b_n^v,
 \qquad
 b_n=\frac{\sqrt{2\pi}\,z_n^{z_n-1/2}e^{-z_n}}{\Gamma(z_n)},
 \qquad \log b_n=-\frac1{6n}+O(n^{-3}).
\]
The saddle densities have forms
\[
 \begin{split}
 f^S_{n,B'}(Q)&=\left(\frac n{2\pi}\right)^{d_v/2}
                 e^{-nI_S(Q)}\mathcal A^S_{B'}(Q,n),\\
 f^G_{n,B'}(Q,0)&=\left(\frac n{2\pi}\right)^{(d_v+v)/2}
                 e^{-nI_G(Q,0)}\mathcal A^G_{B'}(Q,n).
 \end{split}
\]
Their amplitudes at zero target and leading order equal one.

Hence
\[
 \rho_{B'}=
 e^{n(I_G-I_S)}\,
       \mathcal A^S_{B'}b_n^v/\mathcal A^G_{B'}.
\]
Every power of \(n/(2\pi)\), and the entire lattice covolume, have
cancelled before expanding.  As a check, the exact spherical Gram density is
\[
 \frac{\Gamma(n/2)^v}
 {\pi^{d_v/2}\prod_{j=0}^{v-1}\Gamma((n-j)/2)}
 (\det Q_{B'})^{(n-v-1)/2},
\]
whose central factor has the same leading value
\((n/(2\pi))^{d_v/2}\).

The sign transform through degree three is
\[
 \Psi_S(t)=\frac12\sum_{a<b}t_{ab}^2
       +\sum_{a<b<c}t_{ab}t_{ac}t_{bc}+O(|t|^4).
\]
Indeed the only nonzero third products are triangles, each with value
one.  Its Legendre transform therefore has quadratic term
\(\tfrac12\sum Q_{ab}^2\) and cubic term
\(-\sum Q_{ab}Q_{ac}Q_{bc}\).  These agree with the expansion of
\(-\tfrac12\log\det Q_{B'}\).  Thus the rate difference starts in
degree four, as in \eqref{prof:eq-label-expansion}.

Here is the label count for all subsequent terms.  Each statistic
coordinate carries at most two labels, including the single-label
Gaussian length coordinates.  A zero-tilt cumulant tensor with \(M\)
coordinate slots therefore involves at most \(2M\) labels, and its
value is unaffected by additional replicas.  Formal inversion of
\(\nabla\Psi(t)=Q\) produces rooted trees with nonlinear vertex valences
\(m\ge3\).  A degree-\(b\) Legendre coefficient satisfies
\(\sum(m-2)=b-2\).  Since \(m\le3(m-2)\), it uses at most
\(3(b-2)\) cumulant slots, hence at most \(6b\) labels.  Setting the
Gaussian diagonal targets to zero preserves this bound.

For the central amplitude, perform saddle Fourier inversion on the
\(n^{-1/2}\) frequency scale against an identity-covariance Gaussian.
This expands the determinant instead of leaving an implicit
factor depending on dimension.  A higher Fourier vertex of degree
\(m\ge3\) carries weight \(n^{-(m-2)/2}\); at order
\(n^{-j/2}\), their total Fourier-slot count is at most \(3j\).
A vertex from the perturbed Hessian has two Fourier slots and at least
one tilt slot.  If the initial Taylor expansion has \(T\) tilt slots,
all its slots therefore number at most \(3j+3T\).

Substitute the
rooted-tree expansion for the saddle, with a total of \(b\) target
leaves.  This adds at most \(3(b-T)\) slots.  There are thus at most
\(3(j+b)\) slots, or \(6(j+b)\) labels.  Gaussian pairings identify
slots and introduce no further labels; unused coordinates integrate
to one.  This counts the determinant, inverse-Hessian contractions,
and every higher correction.  Odd \(j\) vanish by Fourier parity,
although allowing them in the bookkeeping is harmless.
The factor \(b_n^v\) obeys the same rule by multiplying its
single-label expansions.  Products and reciprocals of series with
leading constant one preserve the bound, because the weights \(b+j\)
add.  This proves the asserted support rule even for coefficients
that depend on \(|B'|\) after their label sums are evaluated.

For fixed \(p\), the transforms are analytic near zero with uniformly
positive Hessian there.  The saddles on
\eqref{prof:eq-bottom-window} lie in this neighborhood.  Taylor
expansion under Gaussian damping on logarithmic frequency balls gives
any required inverse-polynomial relative error.  The sign Fourier
tails are negligible by the reciprocal-lattice contraction just
proved; the Gaussian quadratic determinant formula controls its
remaining Fourier tails.  This proves the uniform saddle remainders.
Moreover \(nQ^4=O(n^{-1/3}\log^{C_p}n)=o(1)\) on the window, so
the exponential in \eqref{prof:eq-label-expansion} can also be
expanded to arbitrary inverse-polynomial precision.

\paragraph{A common integration measure and high-order quadrature.}
Fix the orbit and let \(\mu_o\) be its normalized spherical Gibbs
law, equivalently the Gaussian law below conditioned on length.
Let \(K(x_1,\ldots,x_p)\) be either \(1\) or the product of
squares displayed above, and write
\(\mathcal U(x)=(x_a^Tx_b/L)_{a<b}\). Define the nonnegative
common kernel density \(G_p\) by
\[
 \int_A G_p(U)\,dU
 =\E_{\mu_o^{\otimes p}}
       [K(x_1,\ldots,x_p)\mathbf1_{\{\mathcal U(x)\in A\}}]
\]
for Borel \(A\) in the interior Gram region. This definition
includes the common normalization \(Z_o(h)^{-p}\). We now bound
the density and its derivatives.

Use \(p\) independent Gaussian representatives
\(x_a=m_o+D^{1/2}\eta_a\), with \(D=(\lambda-J)^{-1}\), and impose
all their lengths exactly.  Fourier-invert the length statistics
\((\|x_a\|^2-n)/L\) and overlaps \(U_{ab}=x_a^Tx_b/L\).
For a symmetric Fourier matrix \(T\), the tilted white covariance is
\((I-2iT\otimes D/L)^{-1}\), up to fixed conventions on off-diagonal
coefficients, and has norm at most one.  The tilted mean is bounded
by \(C_p\|T\|\|D^{1/2}m_o\|/L\).

The polylogarithmic edge inputs give
\[
 \|D/L\|+\operatorname{Tr}((D/L)^2)
              +\|D^{1/2}m_o\|/L\le\log^C n,
\]
with logarithmically bounded reciprocals of the exact length densities
in these units.  The spectral spread supplies arbitrarily many
covariance eigenvalues of inverse-polylogarithmic size.  As a
nonzero \(T\) has an eigenvalue at least \(\|T\|/\sqrt p\), the
quadratic determinant is integrable against every fixed frequency
power with only a logarithmic loss; the noncentral characteristic
factor has modulus at most one.

The matrix-moment kernel is a product of squares of
\(\eta_a^T(D/L)\eta_b\) with \(a\ne b\).  Its Wick expansion under
this complex tilt also costs only powers of \(\log n\) and \(T\).
A contraction loop of length at least two is bounded by powers of
\(\|D/L\|\) times \(\operatorname{Tr}((D/L)^2)\).  A loop of length
one uses an off-diagonal replica block of the tilted covariance,
whose identity term vanishes; the resolvent identity supplies an
additional \(D/L\), giving the same bound with a frequency factor.

Open contractions end in the already bounded means.  Differentiation
of the resulting density adds only frequency powers. Consequently
\(G_p(U)\) has all
fixed derivatives bounded by powers of \(\log n\); on the common
smooth window their \(L^1\) norms have the same bounds.  This holds
for both the mass and matrix kernels.

The density and all required derivatives are now controlled. This is
what allows high-order quadrature on each subset's actual measure.
For the subset \(B'\), only its internal coordinates initially lie on
a lattice.  In \(U\) units that lattice has covolume
\(c_v/L^{d_v}\) and a fixed-shape basis of size \(O_p(L^{-1})\).
Its mixed lattice--continuous integral is exactly the quadrature of
\(\rho_{B'}((L/n)U,n)G_p(U)\), with that cell-volume factor.
For every finite smooth expansion term, repeated quadrature after
the fixed basis change, or Poisson summation, gives error
\(O_{p,A}(L^{-A}\log^C n)\), uniformly over affine lattice shifts.

The relative expansion remainder also has a uniform absolute bound
on the common window. Indeed the rate transforms agree through
degree three, and their analytic amplitudes are one at zero, so
for fixed \(p\)
\[
 \rho_{B'}(Q,n)
 =\exp\{O_p(n\max_{a<b}|Q_{ab}|^4)\}
       [1+O_p(\max_{a<b}|Q_{ab}|+n^{-1})]=1+o(1).
\]
Here \(\max|Q_{ab}|=O(n^{-1/3}\log^{C_p}n)\), so the
exponent tends to zero uniformly over all subsets \(B'\).
Thus a relative remainder \(O_{p,A}(n^{-A})\) is also an
absolute remainder of that order. It requires no derivatives, but
is still evaluated on its own mixed lattice--continuous measure.
Apply the same high-order quadrature to the smooth windowed
\(G_p\) itself. Its derivative bounds show that its mixed sum is
at most \(\log^C n+O_{p,A}(L^{-A}\log^C n)\), uniformly in the
affine shift, because its continuous mass has that bound.
Multiplying this mixed-sum bound by the uniform absolute remainder
controls the remainder on its actual measure. We may therefore convert all subset terms to one continuous
integral, with arbitrarily high inverse-polynomial error and with
all covolumes included.  The different subset lattices need not have
matching supports before conversion.

\paragraph{Cancellation and the probability bound.}

Quadrature has put every subset term on one continuous integration
measure. We can now take the alternating sum without identifying the
original lattice supports. A surviving term must use every replica
label, and the coefficient-count rule turns that requirement into a
power of the dimension that grows with the moment order.
A term in the fully expanded expression has rate factors of degrees
\(b_1,\ldots,b_t\ge4\), and an amplitude of degree \(b\) and order
\(n^{-j/2}\).  Put \(W=b_1+\cdots+b_t+b+j\).

Its coefficient involves at most \(6W\) labels.  On the window,
\[
 nQ^{b_i}=O(n^{-b_i/12}\log^{C_p}n),\qquad
 n^{-j/2}Q^b=O(n^{-b/3-j/2}\log^{C_p}n),
\]
where the first estimate uses \(b_i/3-1\ge b_i/12\) for
\(b_i\ge4\).  The whole term thus costs at least
\(n^{-W/12}\), up to logarithms.  A diagram whose label set
\(S\) is a proper subset of \([p]\) cancels exactly, since
\[
 \sum_{B':\,S\subseteq B'\subseteq[p]}(-1)^{p-|B'|}=0.
\]
Every surviving term has \(6W\ge p\), and is bounded by
\(n^{-p/72}\) before logarithmic losses.

Choose the expansion
order sufficiently large in terms of \(p\) and the desired error
power.  The logarithmic kernel bounds and integration, followed by
an enlargement of the threshold in \(n\), give the safe absolute
bound \(n^{-cp}\), for example with \(c=1/200\).

We now return from the expansion to probability. The calculation bounds
the Haar expectation of the truncated alternating power.
The separate second-moment bounds remove the exponentially exceptional
orientations even though the truncated sum may be signed there.  On their complement it agrees with the nonnegative
quantity \(\|\Delta\|_{\mathrm{HS}}^p\) up to negligible error.
For the mass choose threshold \(n^{-c/2}\); Markov then gives Haar
failure at most \(n^{-cp/2}\). For the matrix choose threshold one
in Hilbert--Schmidt norm; its Haar failure is at most \(n^{-cp}\).
The planted likelihood \(R=Z(h)/Z_o(h)\) has bounded second
moment by the two-replica comparison and pair tails.
Cauchy--Schwarz therefore leaves planted failure exponents at
least \(cp/4\) and \(cp/2\), respectively. A sufficiently large
fixed even \(p\) gives any requested \(B\), with room for the
earlier width exclusions.

Let \(M_{\rm cube}\) and \(M_{\rm sphere}\) be the normalized raw
second moments about \(m_o\). On the retained event,
\[
 R=1+O(n^{-c/2}),\qquad
 R M_{\rm cube}-M_{\rm sphere}=L\Delta_{\rm matrix},
 \qquad \|\Delta_{\rm matrix}\|\le
             \|\Delta_{\rm matrix}\|_{\rm HS}\le1.
\]
The spherical bound is
\(\|M_{\rm sphere}\|\le CL\ell^C\). Thus
\(\|M_{\rm cube}\|\le CL\ell^C\), and
\(\Sigma_s\preceq M_{\rm cube}\). This proves \eqref{prof:eq25}.
\end{proof}
